% ============================================================
% f2015ahw_q19_pyramid.tex
% 2015 安徽卷（文科）第 19 题 · 三棱锥 P-ABC —— 清单 F05 / §5
%
% 来源：.work/phase3-b12/figure-inventory.md §5（源 PDF 内容流逐 path 实测真页 pt）
%       源件 .work/phase3-b12/source/2015安徽文.pdf p2
%
% 坐标换算：画布 [x=1pt, y=1pt]（单位 TeX pt）。局部坐标 = 真页 pt − 原点
%   O = (222.468, 240.178)（清单 §5 几何 bbox 左下角）。真页 pt → 局部 pt 为
%   1:1 平移、不缩放——清单 §0.3 的 0.999622 是 pdftocairo 渲染件的 SVG 观测口径换算，
%   只用于与渲染件对读，不入绘图坐标。渲染层面 1 TeX pt = 0.996264 PDF pt（72/72.27），
%   整图统一 0.37%，属排版口径而非坐标换算误差。
%
% 线宽（清单 §0.4 本卷实测分档）：6 条棱一律 0.3985pt（p2 档）
% 虚线：[2.98883 2.98883] 相位 0（仅 AC 棱）
% 箭头：无（本图不含箭头）
% 标签：字号 absolute 10.5pt（顶点字母 CMMI 10.5 数学斜体）；一律 anchor=base 落在清单实测基线上——
%   清单给的是字框中心：汉字字框 = 全角 em 方框（基线 = 中心 − 0.36em），拉丁/数学字框
%   = 字体 ascent/descent 方框（基线 = 中心 − 0.25em）；两式均与清单另给的基线自洽。
%
% 照源复刻的存疑项：
%    - AC 为唯一虚线棱（源 [2.98883 2.98883] 0，与 p1 的 2.989 为同一量的不同舍入，B-S7）；
%      其余五棱实线，线宽 0.3985pt（p2 独有档位，比 p1 的 0.398 多 0.0005pt），照源不复核。
%    - 题面含 ∠BAC = 60°，但源图内无任何角度标注（无角弧、无度数文字、无直角符号、
%      无 BM 线段与 M 点，B-S6）——照源不补标注。
%    - 源每条棱的路径都含 1 个零长度重复子路径（x y m x y m … l S，B-S5），不落墨，已丢弃。
% ============================================================
\begin{tikzpicture}[x=1pt, y=1pt,
  every node/.style={inner sep=0, font=\fontsize{10.5pt}{12.6pt}\selectfont},
  line cap=butt, line join=miter, line width=0.3985pt,
  hidden/.style={dash pattern=on 2.98883pt off 2.98883pt}]

  % ---- 棱（清单 §5.1；PA 竖直、BC 水平）：PA / PB / PC / AB / BC 实线，AC 虚线 ----
  \draw (10.911,113.385) -- (10.911,37.412);
  \draw (10.911,113.385) -- (37.767,10.556);
  \draw (10.911,113.385) -- (169.357,10.556);
  \draw (10.911,37.412) -- (37.767,10.556);
  \draw (37.767,10.556) -- (169.357,10.556);
  \draw[hidden] (10.911,37.412) -- (169.357,10.556);

  % ---- 标签（清单 §5.2；anchor=base 于实测基线）----
  \node[anchor=base] at (10.182,116.767) {$P$};
  \node[anchor=base] at (3.942,33.827) {$A$};
  \node[anchor=base] at (32.572,2.037) {$B$};
  \node[anchor=base] at (173.912,2.037) {$C$};
\end{tikzpicture}
